\documentclass[11pt,a4paper]{beamer}
\usetheme{Madrid}
\usepackage[utf8]{inputenc}
\usepackage{amsmath}
\usepackage{amsfonts}
\usepackage{amssymb}
\usepackage{graphicx}
\usepackage{tikz}
\usepackage{booktabs}
\usepackage{array}
\usepackage{colortbl}
\usepackage{multicol}

% ============================================================
% PROFESSIONAL FONT SETTINGS (REDUCED SIZE ONLY)
% ============================================================
\usefonttheme{professionalfonts}

\setbeamerfont{normal text}{size=\tiny}
\setbeamerfont{itemize/enumerate body}{size=\tiny}
\setbeamerfont{itemize/enumerate subbody}{size=\tiny}
\setbeamerfont{block body}{size=\tiny}
\setbeamerfont{block title}{size=\scriptsize}
\setbeamerfont{frametitle}{size=\small}
\setbeamerfont{title}{size=\small}
\setbeamerfont{subtitle}{size=\footnotesize}
\setbeamerfont{author}{size=\tiny}
\setbeamerfont{date}{size=\tiny}

\setlength{\leftmargini}{0.3em}
\setlength{\leftmarginii}{0.3em}
\setlength{\labelsep}{0.1em}
\setlength{\itemsep}{0.02em}
\setlength{\parskip}{0.02em}
\setlength{\topsep}{0em}
\setlength{\partopsep}{0em}

\newcommand{\redflag}[1]{\textcolor{red}{\textbf{[!] #1}}}
\newcommand{\greennote}[1]{\textcolor{green!50!black}{\textbf{[CORRECT] #1}}}
\newcommand{\examfocus}[1]{\textcolor{orange!80!black}{\textbf{[FOCUS] #1}}}
\newcommand{\trap}[1]{\textcolor{purple}{\textbf{[TRAP] #1}}}
\newcommand{\correct}[1]{\textcolor{green!50!black}{\textbf{[right] #1}}}
\newcommand{\scenario}[1]{\textcolor{blue!70!black}{\textbf{[SCENARIO] #1}}}
\newcommand{\keyterm}[1]{\textcolor{red!80!black}{\textbf{[KEY] #1}}}

\title[CAMS Exam Preparation]{CAMS Exam Preparation}
\subtitle{High Net Worth Individuals - Hard Questions}
\author{Private Banking | Wealth Management | Trusts | PEPs}
\date{}

\setbeamercolor{title}{fg=white,bg=blue!70!black}
\setbeamercolor{frametitle}{fg=white,bg=blue!70!black}
\setbeamertemplate{navigation symbols}{}

\begin{document}
	
	% TITLE SLIDE
	\begin{frame}
		\titlepage
		\centering
		\tiny 30 Hard Questions | Private Banking | Trusts | Offshore Structures | PEPs
	\end{frame}
	
	% ============================================================
	% QUESTION 1 - PRIVATE BANKING CONFLICT OF INTEREST
	% ============================================================
	
	\begin{frame}{Private Banking: Relationship Manager Conflict}
		\tiny
		\scenario{A private bank relationship manager is compensated primarily based on Assets Under Management (AUM) and revenue generated from client accounts. The RM has a client who is a foreign PEP from a high-risk jurisdiction. The client wants to transfer \$15 million through a complex structure involving three shell companies in the Cayman Islands and a trust in the Bahamas. The RM pushes to onboard the client quickly without enhanced due diligence, telling compliance: "This client will generate \$500,000 in annual fees. Don't delay this with unnecessary questions." The compliance officer notices that the RM has already exceeded his annual bonus target and this client would add an additional 30 percent.}
		
		\vspace{0.02cm}
		\textbf{What is the primary inherent financial crime risk in private banking illustrated by this scenario?}
		\vspace{0.02cm}
		
		\begin{itemize}
			\item A) Lack of customer identification technology
			\item B) Conflict of interest where RM prioritizes revenue over AML compliance
			\item C) Insufficient transaction monitoring systems
			\item D) Lack of sanctions screening
		\end{itemize}
		
		\vspace{0.02cm}
		\trap{The problem is technology or process failure.}
		
		\vspace{0.02cm}
		\correct{Answer: B - Private banking RMs compensated based on AUM face inherent conflict of interest. The desire to maintain close relationships and generate fees may cause RMs to overlook warning signs.}
		
		\vspace{0.02cm}
		\redflag{Performance compensation structures in private banking create inherent ML risk.}
	\end{frame}
	
	% ============================================================
	% QUESTION 2 - COMPLEX OWNERSHIP STRUCTURES
	% ============================================================
	
	\begin{frame}{Private Banking: Complex Ownership Structures}
		\tiny
		\scenario{A high net worth individual approaches a private bank with the following structure: The client is the settlor of a trust in the Cook Islands. The trust owns 100 percent of a BVI holding company. The BVI company owns 60 percent of a Cyprus investment fund. The Cyprus fund owns a Swiss bank account with \$25 million. The client refuses to provide any documentation about the trust beneficiaries or the source of the original \$25 million. The client states: "My privacy is protected by these jurisdictions. You don't need to know more." The private bank's relationship manager wants to onboard the client because the client promises to bring an additional \$50 million in assets within 6 months.}
		
		\vspace{0.02cm}
		\textbf{What are the primary AML risks presented by this client structure?}
		\vspace{0.02cm}
		
		\begin{itemize}
			\item A) No risks; trusts are legitimate wealth management tools
			\item B) Complex ownership structure obscures ultimate beneficial owner; multiple offshore jurisdictions with secrecy laws; client refusal to disclose source of funds
			\item C) Only the Swiss bank account is a risk
			\item D) Only the BVI company is a risk
		\end{itemize}
		
		\vspace{0.02cm}
		\trap{Trusts and offshore companies are standard for HNW clients; no additional risk.}
		
		\vspace{0.02cm}
		\correct{Answer: B - Complex multi-jurisdictional structures with opaque beneficial ownership are high risk. Trusts can be the last layer of secrecy. Red flags: multiple offshore jurisdictions, client refusal to provide documentation, pressure to onboard quickly.}
		
		\vspace{0.02cm}
		\redflag{Trusts + offshore companies + secrecy jurisdictions + refusal to disclose = maximum risk.}
	\end{frame}
	
	% ============================================================
	% QUESTION 3 - SOURCE OF WEALTH VS SOURCE OF FUNDS
	% ============================================================
	
	\begin{frame}{Private Banking: Source of Wealth vs Source of Funds}
		\tiny
		\scenario{A private bank is onboarding a new client with \$30 million to invest. The client provides the following information: Source of Wealth: Inheritance from family real estate business sold 10 years ago. Source of Funds for this specific deposit: Sale of cryptocurrency holdings. The client provides a will showing inheritance of \$25 million from his father. However, the bank cannot verify the cryptocurrency sale because the client used a privacy coin (Monero) and an unregulated offshore exchange. The client's father's real estate business was previously investigated for money laundering 15 years ago, but no charges were filed. The client refuses to provide additional documentation about the crypto transaction.}
		
		\vspace{0.02cm}
		\textbf{What is the distinction between Source of Wealth and Source of Funds in this scenario, and what is the bank's obligation?}
		\vspace{0.02cm}
		
		\begin{itemize}
			\item A) Source of Wealth and Source of Funds are the same; no distinction
			\item B) Source of Wealth is the inheritance (\$25 million); Source of Funds is the crypto sale (\$30 million). Bank must verify BOTH independently. The unverifiable crypto sale requires EDD and potential SAR
			\item C) Only Source of Wealth matters for HNW clients
			\item D) Only Source of Funds matters for this deposit
		\end{itemize}
		
		\vspace{0.02cm}
		\trap{Verifying Source of Wealth is sufficient for HNW clients.}
		
		\vspace{0.02cm}
		\correct{Answer: B - Source of Wealth is how client accumulated total wealth over time. Source of Funds is the specific origin of money for this transaction. Both require verification. Unverifiable crypto sale through privacy coin and unregulated exchange is a red flag requiring EDD and potential SAR.}
		
		\vspace{0.02cm}
		\redflag{Source of Wealth and Source of Funds are different. Both must be verified.}
	\end{frame}
	
	% ============================================================
	% QUESTION 4 - TRUST STRUCTURE RED FLAGS
	% ============================================================
	
	\begin{frame}{Private Banking: Trust Structure Red Flags}
		\tiny
		\scenario{A private bank is reviewing an existing trust account established 8 years ago with \$12 million. The trust structure: Settlor: Mr. Smith (deceased 5 years ago). Trustee: Smith Family Office Ltd (a company owned 100 percent by Mr. Smith's son). Beneficiaries: Mr. Smith's three children. The trust has the following activity in the past 12 months: (1) Wire transfer of \$2 million to a company owned by the son's spouse in a high-risk jurisdiction, (2) Purchase of a \$3 million yacht titled in the trust's name, (3) Loan of \$1.5 million to the son's personal account with no repayment terms, (4) Cryptocurrency purchase of \$500,000 through an unregulated VASP. The bank's KYC files show the settlor was a foreign PEP (former minister of finance) but the bank never applied EDD because "the client died." The son/trustee refuses to explain the purpose of the wire transfer or the yacht purchase.}
		
		\vspace{0.02cm}
		\textbf{What red flags indicate potential misuse of this trust for money laundering?}
		\vspace{0.02cm}
		
		\begin{itemize}
			\item A) No red flags; trusts can make any lawful transactions
			\item B) PEP settlor with no EDD; self-dealing transactions (son as trustee benefiting himself); unexplained transfers to high-risk jurisdiction; crypto purchases through unregulated VASP; trustee refusal to explain
			\item C) Only the yacht purchase is a red flag
			\item D) Only the cryptocurrency purchase is a red flag
		\end{itemize}
		
		\vspace{0.02cm}
		\trap{The settlor is deceased; PEP status no longer matters.}
		
		\vspace{0.02cm}
		\correct{Answer: B - PEP status applies to the individual regardless of death when concerns exist about source of funds. Self-dealing (trustee benefiting himself) is a red flag. The bank should have applied EDD to the foreign PEP settlor. Unexplained transactions and refusal to explain require SAR.}
		
		\vspace{0.02cm}
		\redflag{Once a PEP, always a PEP for risk assessment purposes. Trust self-dealing is a red flag.}
	\end{frame}
	
	% ============================================================
	% QUESTION 5 - OFFSHORE FINANCIAL CENTER RISKS
	% ============================================================
	
	\begin{frame}{Private Banking: Offshore Financial Center Risks}
		\tiny
		\scenario{A private bank client maintains accounts in three offshore financial centers: Cayman Islands, British Virgin Islands, and Bermuda. The client is a resident of Country X (not sanctioned). The client's business operates entirely in Country X. The client has no employees, offices, or operations in any offshore jurisdiction. The client's total assets in offshore accounts: \$45 million. The client explains: "I use offshore centers for tax efficiency and asset protection." The bank notices: (1) Round-tripping transactions where funds move from Cayman to BVI to Bermuda and back to Cayman within 30 days, (2) Rapid transfers between the client's three offshore accounts with no apparent economic purpose, (3) The client recently added a shell company in Panama as an additional account signatory, (4) The client is a domestic PEP in Country X (member of parliament).}
		
		\vspace{0.02cm}
		\textbf{Which red flags associated with offshore financial centers are present in this scenario?}
		\vspace{0.02cm}
		
		\begin{itemize}
			\item A) No red flags; OFCs are legitimate for tax planning
			\item B) Complex ownership structures, round tripping, rapid asset transfers between offshore entities, use of shell companies, presence of a PEP
			\item C) Only the PEP status is a red flag
			\item D) Only the round tripping is a red flag
		\end{itemize}
		
		\vspace{0.02cm}
		\trap{Tax planning using OFCs is legitimate; no red flags.}
		
		\vspace{0.02cm}
		\correct{Answer: B - Red flags include: complex structures, round tripping (funds moving in and out with no economic purpose), rapid asset transfers between offshore entities, use of shell companies, and PEP involvement. Each red flag requires enhanced due diligence.}
		
		\vspace{0.02cm}
		\redflag{Round tripping and rapid transfers between OFCs without economic purpose = suspicious.}
	\end{frame}
	
	% ============================================================
	% QUESTION 6 - PRIVATE INVESTMENT COMPANIES
	% ============================================================
	
	\begin{frame}{Private Banking: Private Investment Companies}
		\tiny
		\scenario{A high net worth client establishes a Private Investment Company (PIC) in the British Virgin Islands. The PIC is owned by a trust for which the client is both settlor and beneficiary. The PIC has a bank account at the private bank with \$20 million. The client instructs the bank to execute trades through the PIC account but refuses to provide any information about the PIC's underlying assets or investment strategy. The client says: "PICs are private. You don't need to know what I invest in. Just execute my trades." The bank's transaction monitoring shows the PIC account receives regular wires from three other offshore companies with no apparent connection to the client. The PIC then wires funds to a cryptocurrency exchange in a jurisdiction with weak AML controls.}
		
		\vspace{0.02cm}
		\textbf{Why are Private Investment Companies considered higher risk in private banking?}
		\vspace{0.02cm}
		
		\begin{itemize}
			\item A) PICs are illegal in most jurisdictions
			\item B) PICs reduce transparency of ultimate beneficial owner and can be used to obscure the origin and destination of funds
			\item C) PICs cannot hold bank accounts
			\item D) PICs are only used for illegal purposes
		\end{itemize}
		
		\vspace{0.02cm}
		\trap{PICs are legitimate investment vehicles; no additional risk.}
		
		\vspace{0.02cm}
		\correct{Answer: B - PICs are often established in offshore jurisdictions with secrecy laws. The opaque nature and complex structures make it difficult to trace fund origin and destination. Client refusal to disclose underlying assets is a red flag. Connections to unrelated offshore entities and crypto exchanges increase risk.}
		
		\vspace{0.02cm}
		\redflag{PICs + refusal to disclose underlying assets + crypto = high risk.}
	\end{frame}
	
	% ============================================================
	% QUESTION 7 - TRUST DUE DILIGENCE REQUIREMENTS
	% ============================================================
	
	\begin{frame}{Private Banking: Trust Due Diligence Requirements}
		\tiny
		\scenario{A private bank is asked to open an account for a trust with the following structure: Settlor: Mr. A (deceased). Trustees: Law Firm LLP (professional trustee). Beneficiaries: "The children and grandchildren of Mr. A" (no names provided). Protector: Mrs. B (Mr. A's widow). The trust has existed for 15 years. The trust deed states that the trustees have "absolute discretion" to distribute income and principal. The current trust assets are \$8 million. The law firm trustee refuses to disclose the names of the beneficiaries, citing "beneficiary privacy" and the fact that the trust is discretionary. The bank's KYC team requests the names of all beneficiaries. The law firm refuses.}
		
		\vspace{0.02cm}
		\textbf{Under FATF standards, what due diligence must the bank perform on this trust?}
		\vspace{0.02cm}
		
		\begin{itemize}
			\item A) No due diligence needed; trusts are private arrangements
			\item B) Bank must identify the settlor, all trustees, the protector (if any), AND all beneficiaries, even if discretionary
			\item C) Only the settlor and trustees need to be identified
			\item D) Only the law firm trustee needs to be identified
		\end{itemize}
		
		\vspace{0.02cm}
		\trap{Trusts are private; beneficiaries do not need to be identified.}
		
		\vspace{0.02cm}
		\correct{Answer: B - Banks must identify settlor, all trustees, protector, and all beneficiaries of a trust. For discretionary trusts where beneficiaries are not yet determined, the bank must understand the class of beneficiaries and obtain information about known beneficiaries. Refusal to provide beneficiary information is grounds to reject the account.}
		
		\vspace{0.02cm}
		\redflag{All trust beneficiaries must be identified, including discretionary beneficiaries.}
	\end{frame}
	
	% ============================================================
	% QUESTION 8 - SPECIAL PURPOSE VEHICLE MISUSE
	% ============================================================
	
	\begin{frame}{Private Banking: Special Purpose Vehicle Risks}
		\tiny
		\scenario{A private banking client establishes a Special Purpose Vehicle (SPV) in Delaware to hold a single asset: a \$12 million painting. The client then uses the SPV as collateral for a \$9 million loan from a different bank. The loan proceeds are wired to the private bank and used to purchase cryptocurrency. The client then creates a second SPV in Luxembourg to hold the cryptocurrency. The client refuses to explain the business purpose of the SPVs beyond "asset protection." The client's wealth manager notes that the client has created 7 SPVs in the past 3 years, each dissolved within 12-18 months after transferring assets.}
		
		\vspace{0.02cm}
		\textbf{What financial crime risks are associated with the client's use of SPVs?}
		\vspace{0.02cm}
		
		\begin{itemize}
			\item A) No risks; SPVs are standard for single assets
			\item B) Complex opaque structures to disguise true beneficial ownership; rapid creation and dissolution of SPVs; use of SPVs to obscure source of funds for cryptocurrency; no clear legitimate business purpose
			\item C) Only the loan transaction is a risk
			\item D) Only the cryptocurrency purchase is a risk
		\end{itemize}
		
		\vspace{0.02cm}
		\trap{SPVs are legitimate for holding assets; no ML risk.}
		
		\vspace{0.02cm}
		\correct{Answer: B - SPVs can have complex opaque structures to disguise true beneficial ownership. Red flags: multiple SPVs created and dissolved rapidly, no clear business purpose, use of SPVs to obscure funds moving into crypto. EDD required.}
		
		\vspace{0.02cm}
		\redflag{Rapid creation and dissolution of SPVs + crypto + no business purpose = red flags.}
	\end{frame}
	
	% ============================================================
	% QUESTION 9 - SOVEREIGN WEALTH FUND RISKS
	% ============================================================
	
	\begin{frame}{Private Banking: Sovereign Wealth Fund Risks}
		\tiny
		\scenario{A private bank is onboarding a Sovereign Wealth Fund (SWF) from a country with high corruption perception index ranking. The SWF has \$50 billion in assets under management. The SWF's board includes the country's finance minister (a foreign PEP), the central bank governor (foreign PEP), and two individuals whose backgrounds cannot be verified. The SWF wants to invest \$200 million in the private bank's funds. The SWF refuses to provide any information about source of funds, stating: "We are a sovereign entity. We are exempt from AML requirements by international law." The bank's compliance officer notes that the SWF's country has been on FATF Grey List for the past 18 months.}
		
		\vspace{0.02cm}
		\textbf{What is the bank's obligation regarding this Sovereign Wealth Fund?}
		\vspace{0.02cm}
		
		\begin{itemize}
			\item A) No obligations; SWFs are exempt from AML
			\item B) SWFs are not automatically exempt; bank must apply EDD, verify beneficial owners (including PEP board members), assess source of funds, and consider Grey List status
			\item C) Only need to verify the SWF's registration
			\item D) Only need to screen board members against sanctions
		\end{itemize}
		
		\vspace{0.02cm}
		\trap{SWFs are government entities exempt from AML.}
		
		\vspace{0.02cm}
		\correct{Answer: B - SWFs are not automatically exempt. EDD required for PEP board members. Source of funds must be verified. FATF Grey List status requires enhanced due diligence. SWF's claim of exemption is incorrect.}
		
		\vspace{0.02cm}
		\redflag{SWFs are not automatically exempt from AML. PEP board members require EDD.}
	\end{frame}
	
	% ============================================================
	% QUESTION 10 - HIGH VALUE ASSET RED FLAGS
	% ============================================================
	
	\begin{frame}{Private Banking: High Value Asset Red Flags}
		\tiny
		\scenario{A private banking client purchases the following high value assets over 6 months: (1) A \$25 million painting at auction paid by wire from an account in the Seychelles, (2) A \$12 million diamond necklace paid in three cash installments of \$4 million each, (3) A \$40 million yacht registered to a shell company in the Marshall Islands, (4) A \$8 million gold bar purchase through a dealer with no AML program. The client's declared annual income from his business is \$3 million. The client cannot explain the source of funds for any purchase. The client's wealth manager notes that the client sells each asset within 12-18 months, often at a loss.}
		
		\vspace{0.02cm}
		\textbf{Which red flags indicate that high value assets are being used for money laundering?}
		\vspace{0.02cm}
		
		\begin{itemize}
			\item A) No red flags; HNW individuals buy luxury assets regularly
			\item B) Funds from high-risk jurisdiction (Seychelles); cash structuring (\$4 million installments); shell company ownership; dealer with no AML; income inconsistency; rapid resale at loss
			\item C) Only the cash payments are a red flag
			\item D) Only the shell company registration is a red flag
		\end{itemize}
		
		\vspace{0.02cm}
		\trap{HNW individuals buying luxury assets is normal.}
		
		\vspace{0.02cm}
		\correct{Answer: B - Red flags: funds from secrecy jurisdiction, cash structuring, shell company ownership, unregulated dealer, income inconsistency ( \$3M income vs \$85M purchases ), rapid resale at loss (criminals willing to accept loss to clean money).}
		
		\vspace{0.02cm}
		\redflag{Rapid resale of luxury assets at a loss = money laundering (the loss is the laundering cost).}
	\end{frame}
	
	% ============================================================
	% QUESTION 11 - FAMILY OFFICE RISKS
	% ============================================================
	
	\begin{frame}{Private Banking: Family Office Risks}
		\tiny
		\scenario{A private bank is opening an account for a Single Family Office (SFO) managing \$500 million for a ultra-high net worth family. The family patriarch was a former government official (PEP) who left office 8 years ago. The family office is registered in Delaware but managed from Switzerland. The family office has the following employees: 2 family members, 3 professional staff, and 1 external investment advisor. The family office refuses to provide information about the external advisor, stating "he is our trusted partner for 20 years." The bank's transaction monitoring shows the family office sends wires to 15 different countries monthly, including two FATF Grey List jurisdictions. The family office also receives regular wires from a company in the British Virgin Islands that is not related to any known family business.}
		
		\vspace{0.02cm}
		\textbf{What are the primary AML risks associated with this Family Office?}
		\vspace{0.02cm}
		
		\begin{itemize}
			\item A) No risks; family offices are low risk
			\item B) PEP involvement (patriarch); opaque structure (Delaware registration, Swiss management); unidentified third parties (external advisor); transactions to Grey List jurisdictions; unexplained incoming wires from BVI
			\item C) Only the PEP status is a risk
			\item D) Only the Grey List transactions are a risk
		\end{itemize}
		
		\vspace{0.02cm}
		\trap{Family offices are legitimate; patriarch left office 8 years ago, no longer PEP.}
		
		\vspace{0.02cm}
		\correct{Answer: B - PEP risk may continue (once a PEP always a PEP). Multi-jurisdictional structure reduces transparency. Unidentified third parties are a red flag. Unexplained incoming wires from shell company jurisdictions require investigation.}
		
		\vspace{0.02cm}
		\redflag{Once a PEP, always a PEP for risk purposes. Unidentified third parties in family offices = red flag.}
	\end{frame}
	
	% ============================================================
	% QUESTION 12 - INSURANCE PRODUCTS FOR HNW
	% ============================================================
	
	\begin{frame}{Private Banking: Insurance Products for HNW Clients}
		\tiny
		\scenario{A private banking client purchases a single-premium variable life insurance policy for \$10 million. The client pays the premium using funds from an offshore account in the Cayman Islands. Three months later, the client takes a policy loan of \$8 million (80 percent of cash value). The loan proceeds are wired to a real estate developer to purchase a luxury condo. Six months later, the client surrenders the policy and receives the remaining cash value of \$1.5 million as a check from the insurance company. The client's wealth manager notes that the client has done this same pattern with three different insurance companies in the past 2 years: purchase policy, take large loan, surrender policy after 6-12 months.}
		
		\vspace{0.02cm}
		\textbf{What money laundering technique is being attempted through these insurance products?}
		\vspace{0.02cm}
		
		\begin{itemize}
			\item A) No laundering; insurance is a legitimate investment
			\item B) Policy loan layering: Illicit premium converted to legitimate loan proceeds; early surrender for clean check; repeating pattern indicates systematic laundering
			\item C) Only the surrender proceeds are laundered
			\item D) Only the loan proceeds are laundered
		\end{itemize}
		
		\vspace{0.02cm}
		\trap{Policy loans are legitimate; no ML risk.}
		
		\vspace{0.02cm}
		\correct{Answer: B - The pattern is classic insurance money laundering: (1) Purchase policy with illicit funds, (2) Take policy loan to receive clean funds (layering), (3) Surrender policy for additional clean funds (integration). Repeating pattern with multiple insurers is a strong red flag.}
		
		\vspace{0.02cm}
		\redflag{Policy loans from insurance policies purchased with illicit funds = layering.}
	\end{frame}
	
	% ============================================================
	% QUESTION 13 - CRYPTO AND HNW CLIENTS
	% ============================================================
	
	\begin{frame}{Private Banking: Cryptocurrency and HNW Clients}
		\tiny
		\scenario{A private banking client with a declared net worth of \$50 million (from real estate development) suddenly begins engaging in significant cryptocurrency activity. Over 6 months: (1) Purchased \$8 million in Bitcoin through a VASP with weak KYC, (2) Transferred Bitcoin to a DeFi protocol and earned 12 percent APY, (3) Converted profits to USDC through a decentralized exchange, (4) Sent USDC to a regulated VASP in Switzerland, (5) Cashed out \$9.5 million to his private bank account. The client cannot explain how he learned about DeFi or why he used multiple platforms. The client's real estate business has shown declining revenues for 3 years. The client refuses to provide wallet addresses or transaction history from the DeFi protocol.}
		
		\vspace{0.02cm}
		\textbf{What red flags are present regarding this HNW client's cryptocurrency activity?}
		\vspace{0.02cm}
		
		\begin{itemize}
			\item A) No red flags; HNW clients invest in crypto
			\item B) Disconnect between business performance and crypto activity; use of VASP with weak KYC; DeFi protocol with no transparency; refusal to provide wallet history; unexplained rapid profit
			\item C) Only the VASP choice is a red flag
			\item D) Only the profit amount is a red flag
		\end{itemize}
		
		\vspace{0.02cm}
		\trap{Crypto investing is normal for HNW clients.}
		
		\vspace{0.02cm}
		\correct{Answer: B - Red flags: declining business revenues but \$8M crypto purchase; unregulated VASP; DeFi opacity; refusal to provide transaction history; unexplained profits. EDD required.}
		
		\vspace{0.02cm}
		\redflag{Declining business income + large crypto purchases = potential source of funds issue.}
	\end{frame}
	
	% ============================================================
	% QUESTION 14 - ART AND COLLECTIBLES LAUNDERING
	% ============================================================
	
	\begin{frame}{Private Banking: Art and Collectibles Laundering}
		\tiny
		\scenario{A private banking client has purchased \$50 million in fine art over 3 years. The bank notices the following pattern: (1) The client buys art through an auction house in London, (2) The art is shipped to a freeport in Geneva, (3) The client never takes possession of the art, (4) The art is sold 6-12 months later to a different buyer, often at a significant loss (20-30 percent), (5) The sale proceeds are deposited to a different bank account than the purchase account. The client has created three separate LLCs, each used for a single art purchase and sale then dissolved. The client refuses to explain the art investment strategy, stating "art is a private matter."}
		
		\vspace{0.02cm}
		\textbf{What indicates that art is being used for money laundering rather than investment?}
		\vspace{0.02cm}
		
		\begin{itemize}
			\item A) No indicators; art investment is legitimate
			\item B) Freeport storage (no physical possession); rapid resale at loss; use of single-purpose LLCs; refusal to explain; different bank accounts for purchase and sale
			\item C) Only the freeport storage is suspicious
			\item D) Only the resale at loss is suspicious
		\end{itemize}
		
		\vspace{0.02cm}
		\trap{Art is a legitimate investment for HNW clients.}
		
		\vspace{0.02cm}
		\correct{Answer: B - Art ML indicators: freeport storage (never physically possessed), rapid resale at loss (criminals accept loss to clean money), single-purpose LLCs (create false provenance), different accounts (layering), refusal to explain.}
		
		\vspace{0.02cm}
		\redflag{Rapid resale of art at a loss + freeport storage = money laundering.}
	\end{frame}
	
	% ============================================================
	% QUESTION 15 - SECURED LOANS FROM PRIVATE BANKING
	% ============================================================
	
	\begin{frame}{Private Banking: Secured Loan Risks}
		\tiny
		\scenario{A private banking client has a \$15 million portfolio of publicly traded stocks and bonds. The client requests a secured loan of \$12 million (80 percent loan-to-value) using the portfolio as collateral. The client states the loan is for "business expansion." The bank approves the loan. Three days after receiving the loan proceeds, the client wires \$10 million to a company in a high-risk jurisdiction with no apparent connection to the client's business. The client then sells the collateral portfolio (triggering capital gains tax) and uses the proceeds to repay the loan 2 months early. The client's business has no record of receiving the \$10 million.}
		
		\vspace{0.02cm}
		\textbf{How could this secured loan structure be used for money laundering?}
		\vspace{0.02cm}
		
		\begin{itemize}
			\item A) Loans cannot be used for money laundering
			\item B) Illicit funds used as collateral; loan proceeds are clean; early repayment after selling collateral creates legitimate-appearing source of funds
			\item C) Only the early repayment is suspicious
			\item D) Only the wire to high-risk jurisdiction is suspicious
		\end{itemize}
		
		\vspace{0.02cm}
		\trap{Secured loans are low risk; collateral guarantees repayment.}
		
		\vspace{0.02cm}
		\correct{Answer: B - Money launderers can use illicit funds as collateral (the portfolio may be purchased with dirty money). The loan proceeds are clean and can be moved freely. Early repayment after selling collateral creates a legitimate paper trail. The wire to high-risk jurisdiction and business with no record of receipt are red flags.}
		
		\vspace{0.02cm}
		\redflag{Secured loans from private banking can launder funds if collateral is purchased with illicit money.}
	\end{frame}
	
	% ============================================================
	% QUESTION 16 - MULTI-JURISDICTION TAX ISSUES
	% ============================================================
	
	\begin{frame}{Private Banking: Multi-Jurisdiction Tax Issues}
		\tiny
		\scenario{A private banking client is a citizen of Country A, resident of Country B for tax purposes, operates business in Country C, and has bank accounts in Countries D, E, and F. The client's declared total assets are \$200 million. The bank notices: (1) The client moves funds between accounts in Countries D, E, and F weekly with no apparent business purpose, (2) The client's tax returns in Country B show only \$2 million in annual income, (3) The client refuses to sign a Common Reporting Standard (CRS) self-certification form, stating "my tax affairs are complex," (4) The client has requested that the bank not report interest income to any tax authority. The bank's jurisdiction is a CRS participant.}
		
		\vspace{0.02cm}
		\textbf{What are the bank's obligations regarding this client's potential tax evasion?}
		\vspace{0.02cm}
		
		\begin{itemize}
			\item A) No obligations; tax evasion is not a predicate offense for money laundering
			\item B) Tax evasion is a predicate offense for money laundering; bank must file CRS report despite client refusal; unexplained cross-border movements and income inconsistency require EDD and potential SAR
			\item C) Only need to file CRS if client signs form
			\item D) Bank cannot report without client consent
		\end{itemize}
		
		\vspace{0.02cm}
		\trap{Tax evasion is not a bank's concern; only money laundering matters.}
		
		\vspace{0.02cm}
		\correct{Answer: B - Tax evasion is a predicate offense for money laundering in most jurisdictions. CRS requires reporting regardless of client consent. Unusual cross-border movements and income inconsistency are red flags. Refusal to provide CRS form is itself a red flag requiring EDD.}
		
		\vspace{0.02cm}
		\redflag{Tax evasion = predicate offense for money laundering. CRS reporting mandatory.}
	\end{frame}
	
	% ============================================================
	% QUESTION 17 - SHELL COMPANY NETWORKS
	% ============================================================
	
	\begin{frame}{Private Banking: Shell Company Networks}
		\tiny
		\scenario{A private banking client controls a network of 25 shell companies across 8 jurisdictions: BVI, Cayman, Delaware, Panama, Seychelles, Bahamas, Belize, and Mauritius. The client is the ultimate beneficial owner of all 25 companies. The client's declared business is international consulting with \$5 million annual revenue. The bank's transaction monitoring shows: (1) Funds move in a circular pattern among all 25 companies, (2) Total annual flow through the network is \$500 million, (3) Each company has a different bank account at a different financial institution, (4) The client refuses to provide financial statements for any company, (5) The client recently added a trust in the Cook Islands as the owner of the entire network, with the client as both settlor and beneficiary.}
		
		\vspace{0.02cm}
		\textbf{What is the primary money laundering risk presented by this shell company network?}
		\vspace{0.02cm}
		
		\begin{itemize}
			\item A) No risk; shell companies are legal business structures
			\item B) Opaque ownership structure obscures fund movement; circular transactions indicate layering; revenue-to-flow ratio ( \$5M revenue vs \$500M flow) is wildly inconsistent; trust overlay adds final secrecy layer
			\item C) Only the circular transactions are a risk
			\item D) Only the revenue inconsistency is a risk
		\end{itemize}
		
		\vspace{0.02cm}
		\trap{Shell companies are legitimate; no risk.}
		
		\vspace{0.02cm}
		\correct{Answer: B - Massive discrepancy between declared revenue (\$5M) and transaction flow (\$500M) is the strongest red flag. Circular transactions indicate layering. Multiple jurisdictions and trust overlay obscure beneficial ownership. This is a sophisticated ML network.}
		
		\vspace{0.02cm}
		\redflag{Revenue vs transaction flow discrepancy = strongest red flag for shell company networks.}
	\end{frame}
	
	% ============================================================
	% QUESTION 18 - ROUND TRIPPING
	% ============================================================
	
	\begin{frame}{Private Banking: Round Tripping Detection}
		\tiny
		\scenario{A private banking client has accounts at two different banks. The pattern over 12 months: (1) Bank A account sends \$1 million to Bank B account, (2) Bank B account sends \$980,000 back to Bank A account 5 days later, (3) Repeat with different amounts: \$2M sent, \$1.95M returned; \$3M sent, \$2.92M returned; \$5M sent, \$4.88M returned. The client's declared business is manufacturing with no cross-border operations. The client cannot explain the purpose of these round-trip transfers. Each round trip loses approximately 2-5 percent of the value (the client's cost). The bank calculates that the client has moved \$50 million through this pattern in 12 months, losing approximately \$2 million in the process.}
		
		\vspace{0.02cm}
		\textbf{What is the most likely purpose of these round-tripping transactions?}
		\vspace{0.02cm}
		
		\begin{itemize}
			\item A) Currency hedging
			\item B) Money laundering layering: creating complex transaction history to obscure fund origin, willing to accept losses (laundering cost)
			\item C) Tax avoidance
			\item D) Legitimate cash management
		\end{itemize}
		
		\vspace{0.02cm}
		\trap{This could be legitimate treasury management.}
		
		\vspace{0.02cm}
		\correct{Answer: B - Round tripping (moving funds out and back) with consistent loss of 2-5 percent indicates layering. Criminals accept losses to clean money. No legitimate business purpose. Inability to explain is a red flag requiring SAR.}
		
		\vspace{0.02cm}
		\redflag{Round tripping with consistent losses = money laundering (loss is the laundering fee).}
	\end{frame}
	
	% ============================================================
	% QUESTION 19 - FOREIGN PEP ENHANCED DUE DILIGENCE
	% ============================================================
	
	\begin{frame}{Private Banking: Foreign PEP Enhanced Due Diligence}
		\tiny
		\scenario{A private bank is onboarding a foreign PEP: a former prime minister of a country that left office 5 years ago. The PEP wants to deposit \$25 million. The PEP provides the following: Source of wealth: Pension from government service ( \$500,000 annually) plus consulting fees from 3 international companies ( \$2 million annually). The PEP cannot explain how \$25 million accumulated from \$2.5 million annual income over 5 years (maximum \$12.5 million before taxes and expenses). The PEP states: "I had savings before becoming prime minister." The PEP refuses to provide tax returns or bank statements from before his government service. The PEP's country is ranked 145 out of 180 on Transparency International's Corruption Perceptions Index.}
		
		\vspace{0.02cm}
		\textbf{Under FATF standards, what enhanced due diligence must the bank perform for this foreign PEP?}
		\vspace{0.02cm}
		
		\begin{itemize}
			\item A) No EDD needed; PEP left office 5 years ago
			\item B) EDD required: verify source of wealth and source of funds; obtain senior management approval; enhanced monitoring; mathematical inconsistency ( \$25M vs \$12.5M max possible) requires SAR if not resolved
			\item C) Only need senior management approval
			\item D) Only need source of funds verification
		\end{itemize}
		
		\vspace{0.02cm}
		\trap{PEP status expires after leaving office.}
		
		\vspace{0.02cm}
		\correct{Answer: B - FATF requires EDD for foreign PEPs regardless of time since leaving office. Bank must verify source of wealth and source of funds. Mathematical inconsistency requires explanation. Refusal to provide documentation and high CPI ranking are additional risk factors. SAR required if discrepancy cannot be resolved.}
		
		\vspace{0.02cm}
		\redflag{Foreign PEPs always require EDD, regardless of when they left office.}
	\end{frame}
	
	% ============================================================
	% QUESTION 20 - DOMESTIC PEP RISK ASSESSMENT
	% ============================================================
	
	\begin{frame}{Private Banking: Domestic PEP Risk Assessment}
		\tiny
		\scenario{A private bank is onboarding a domestic PEP: a senior government official in the same country where the bank operates. The official's declared income is \$300,000 annually from government salary. The official wants to deposit \$2 million from "family inheritance." The bank verifies that the official's father died 2 years ago and left an estate valued at \$1.5 million (after taxes). The additional \$500,000 cannot be explained. The official's son has a company that won 3 government contracts worth \$800,000 in the past year. The official was involved in awarding those contracts. The bank's country has strong anti-corruption laws but weak enforcement.}
		
		\vspace{0.02cm}
		\textbf{Under FATF standards, how should the bank assess this domestic PEP's risk?}
		\vspace{0.02cm}
		
		\begin{itemize}
			\item A) Domestic PEPs are always low risk
			\item B) FATF recommends risk-based assessment for domestic PEPs; red flags include unexplained wealth ( \$500,000 gap) and son's government contracts with father's involvement
			\item C) Domestic PEPs require same EDD as foreign PEPs
			\item D) No EDD needed for domestic PEPs
		\end{itemize}
		
		\vspace{0.02cm}
		\trap{Domestic PEPs are exempt from EDD under FATF.}
		
		\vspace{0.02cm}
		\correct{Answer: B - FATF recommends risk-based approach for domestic PEPs, not automatic EDD. However, red flags in this case (unexplained wealth, son's government contracts) require EDD regardless of PEP type. The \$500,000 unexplained gap and potential conflict of interest are significant.}
		
		\vspace{0.02cm}
		\redflag{Domestic PEPs require risk-based assessment. Unexplained wealth requires EDD.}
	\end{frame}
	
	% ============================================================
	% QUESTION 21 - INTERNATIONAL ORGANIZATION PEP
	% ============================================================
	
	\begin{frame}{Private Banking: International Organization PEP}
		\tiny
		\scenario{A private bank is onboarding a client who is a senior executive at the World Bank (an international organization). The client's declared annual income is \$400,000. The client wants to deposit \$3 million from "savings and investments over 20 years." The client provides bank statements showing monthly deposits of \$10,000-15,000 into a savings account for 20 years, totaling approximately \$3 million. However, the bank notices that the deposits stop during the client's 5 years of employment at the World Bank and resume after leaving. The client cannot explain why deposits stopped during World Bank employment. The client's country of origin is a high-risk jurisdiction.}
		
		\vspace{0.02cm}
		\textbf{What is the bank's obligation regarding this International Organization PEP?}
		\vspace{0.02cm}
		
		\begin{itemize}
			\item A) No EDD needed; international organization PEPs are low risk
			\item B) FATF recommends risk-based assessment; deposit pattern inconsistency (no deposits during World Bank employment) is a red flag requiring EDD
			\item C) Same as foreign PEP: automatic EDD required
			\item D) No obligations; World Bank has its own AML
		\end{itemize}
		
		\vspace{0.02cm}
		\trap{International organization PEPs are exempt from EDD.}
		
		\vspace{0.02cm}
		\correct{Answer: B - FATF recommends risk-based assessment for international organization PEPs. The inconsistency in deposit pattern (deposits stopped during World Bank employment) is a red flag. The client cannot explain why. High-risk country of origin is additional factor. EDD required.}
		
		\vspace{0.02cm}
		\redflag{International organization PEPs require risk-based assessment. Inconsistent wealth accumulation patterns require EDD.}
	\end{frame}
	
	% ============================================================
	% QUESTION 22 - IMMEDIATE FAMILY OF PEP
	% ============================================================
	
	\begin{frame}{Private Banking: Immediate Family of PEP}
		\tiny
		\scenario{A private bank is onboarding a client who is the spouse of a foreign PEP (the spouse is a current cabinet minister in Country X). The spouse has her own business: a consulting firm with \$500,000 annual revenue. The spouse wants to deposit \$10 million into a private banking account. The source of funds is described as "gift from husband." The husband's declared government salary is \$150,000 annually. The husband has no other declared source of wealth. The family's lifestyle includes a \$5 million home, private jets, and luxury vacations. The spouse refuses to provide any documentation about the gift, stating "my husband's finances are private."}
		
		\vspace{0.02cm}
		\textbf{How does FATF treat immediate family of PEPs, and what is the bank's obligation?}
		\vspace{0.02cm}
		
		\begin{itemize}
			\item A) Family members are not PEPs; no EDD needed
			\item B) FATF definition includes immediate family; EDD required; \$10M gift from \$150k salary is mathematically impossible; SAR required
			\item C) Only the PEP spouse requires EDD, not the family member
			\item D) No EDD needed for gifts
		\end{itemize}
		
		\vspace{0.02cm}
		\trap{PEP definition applies only to the individual, not family.}
		
		\vspace{0.02cm}
		\correct{Answer: B - FATF definition of PEP includes immediate family members. The spouse is a PEP by association. The gift amount (\$10M) is wildly inconsistent with government salary (\$150k). Impossible to accumulate from legitimate income. SAR required.}
		
		\vspace{0.02cm}
		\redflag{Immediate family of PEPs are also PEPs. Gifts from PEPs require source of funds verification.}
	\end{frame}
	
	% ============================================================
	% QUESTION 23 - CLOSE ASSOCIATES OF PEP
	% ============================================================
	
	\begin{frame}{Private Banking: Close Associates of PEP}
		\tiny
		\scenario{A private bank is onboarding a client who is a long-time friend and business partner of a foreign PEP (the PEP is a former intelligence agency director). The client and the PEP co-own a company that has won 15 government contracts worth \$200 million over 8 years. The client is the public-facing owner; the PEP's ownership is hidden through a trust. The client's declared personal net worth is \$50 million. The client refuses to identify the other owners of the company, stating "business partners are private." The bank's open source research shows news articles alleging that the PEP used his position to steer contracts to the company. No charges have been filed.}
		
		\vspace{0.02cm}
		\textbf{Under FATF standards, how should the bank treat this close associate of a PEP?}
		\vspace{0.02cm}
		
		\begin{itemize}
			\item A) No special treatment; only the PEP matters
			\item B) Close associates are included in PEP definition; hidden ownership structure; allegations of contract steering; refusal to identify owners; all require EDD and potential SAR
			\item C) Only need to screen the client against PEP lists
			\item D) No action unless charges are filed
		\end{itemize}
		
		\vspace{0.02cm}
		\trap{Close associates are not PEPs; only the individual is.}
		
		\vspace{0.02cm}
		\correct{Answer: B - FATF definition includes close associates of PEPs. Hidden ownership through trust, refusal to identify beneficial owners, and allegations of corruption are red flags. EDD required. Bank should consider SAR.}
		
		\vspace{0.02cm}
		\redflag{Close associates of PEPs are also considered PEPs for AML purposes.}
	\end{frame}
	
	% ============================================================
	% QUESTION 24 - CORRESPONDENT BANKING FOR PRIVATE BANK
	% ============================================================
	
	\begin{frame}{Private Banking: Correspondent Banking Risks}
		\tiny
		\scenario{A private bank in Country A has a correspondent banking relationship with a private bank in Country B (a known offshore financial center). The Country B bank has 500 private banking clients, including 50 foreign PEPs. The Country B bank refuses to provide the names of any clients, citing bank secrecy laws. The Country B bank processes approximately \$10 billion annually in wire transfers for its private banking clients. The Country B bank has been cited by its local regulator for AML deficiencies twice in the past 3 years. The Country A bank's due diligence questionnaire reveals that the Country B bank has no independent AML audit in 4 years.}
		
		\vspace{0.02cm}
		\textbf{What is the Country A bank's obligation under FATF Recommendation 13 (Correspondent Banking)?}
		\vspace{0.02cm}
		
		\begin{itemize}
			\item A) No obligation; correspondent banking is standard
			\item B) Must terminate or restrict relationship; respondent bank has AML deficiencies, refuses to identify underlying clients (including PEPs), operates in offshore center, no independent audit
			\item C) Only need to file annual report
			\item D) Only need to monitor transactions
		\end{itemize}
		
		\vspace{0.02cm}
		\trap{Correspondent banking risks can be managed with monitoring.}
		
		\vspace{0.02cm}
		\correct{Answer: B - FATF Rec. 13 requires EDD for correspondent relationships. Red flags: respondent refuses to identify underlying clients (including PEPs), AML deficiencies, no independent audit, offshore jurisdiction. Termination or significant restriction is appropriate.}
		
		\vspace{0.02cm}
		\redflag{Correspondent bank that refuses to identify underlying private banking clients should be terminated.}
	\end{frame}
	
	% ============================================================
	% QUESTION 25 - NESTED ACCOUNTS IN PRIVATE BANKING
	% ============================================================
	
	\begin{frame}{Private Banking: Nested Account Risks}
		\tiny
		\scenario{A private bank has a correspondent account for Respondent Bank R. Bank R allows its own private banking clients to use the correspondent account directly (nested accounts). The private bank discovers that one of Bank R's private banking clients is a foreign PEP from a sanctioned country. The private bank never approved this nested relationship. Bank R refuses to terminate the PEP client's access, stating "the client is important to our business." The private bank's investigation reveals that 30 percent of the transactions through its correspondent account are for Bank R's private banking clients, not for Bank R itself. The private bank has no KYC information for any of these underlying clients.}
		
		\vspace{0.02cm}
		\textbf{What action should the private bank take regarding this unauthorized nested account arrangement?}
		\vspace{0.02cm}
		
		\begin{itemize}
			\item A) No action; nesting is standard industry practice
			\item B) Unauthorized nesting is a serious risk; private bank must file SAR, demand termination of nesting, and consider terminating correspondent relationship
			\item C) Only need to file SAR for the PEP client
			\item D) Only need to monitor transactions more closely
		\end{itemize}
		
		\vspace{0.02cm}
		\trap{Nesting is permitted if respondent bank has adequate controls.}
		
		\vspace{0.02cm}
		\correct{Answer: B - Unauthorized nesting is a major risk. The correspondent bank has no visibility into underlying clients including PEPs from sanctioned countries. Respondent bank's refusal to terminate nesting is unacceptable. SAR required. Termination of correspondent relationship should be considered.}
		
		\vspace{0.02cm}
		\redflag{Unauthorized nesting of private banking clients through correspondent accounts = terminate relationship.}
	\end{frame}
	
	% ============================================================
	% QUESTION 26 - PRIVATE BANKING EXIT STRATEGY
	% ============================================================
	
	\begin{frame}{Private Banking: Client Exit Strategy}
		\tiny
		\scenario{A private bank has decided to exit a high-risk client relationship: a foreign PEP with unexplained wealth of \$30 million who refuses to provide source of funds documentation. The client has outstanding loans of \$8 million secured by real estate. The client also has a \$10 million irrevocable letter of credit issued by the bank on behalf of the client's business. The client threatens to sue the bank for breach of contract if the accounts are closed. The bank's legal team advises that closing the accounts could trigger litigation. The relationship manager urges the bank to keep the client, citing \$500,000 in annual fees.}
		
		\vspace{0.02cm}
		\textbf{What is the bank's obligation regarding exiting this high-risk private banking client?}
		\vspace{0.02cm}
		
		\begin{itemize}
			\item A) Keep the client; cannot exit due to loans and LCs
			\item B) Exit strategy must consider legal obligations but cannot retain client solely for revenue; manage loan and LC separately; file SAR; document exit decision
			\item C) Only exit after loans are repaid
			\item D) Never exit PEP clients; manage with EDD only
		\end{itemize}
		
		\vspace{0.02cm}
		\trap{Banks cannot exit clients with outstanding credit facilities.}
		
		\vspace{0.02cm}
		\correct{Answer: B - Banks can and should exit high-risk clients even with outstanding facilities. Exit must be managed carefully to avoid breach of contract. Loans and LCs can be allowed to run off or called if contract permits. Revenue cannot be the reason to retain a high-risk client. SAR must be filed.}
		
		\vspace{0.02cm}
		\redflag{Revenue and outstanding loans are not valid reasons to retain a high-risk client.}
	\end{frame}
	
	% ============================================================
	% QUESTION 27 - PRIVATE BANKING CULTURE OF COMPLIANCE
	% ============================================================
	
	\begin{frame}{Private Banking: Culture of Compliance}
		\tiny
		\scenario{A private bank's internal audit reveals the following: (1) The bank's 5 largest clients (all foreign PEPs) contribute 40 percent of the private banking division's revenue, (2) The relationship managers for these 5 clients have never filed a SAR despite significant red flags, (3) The bank's compliance department has recommended terminating 3 of the 5 clients, (4) Senior management has rejected all 3 termination recommendations citing "commercial sensitivity," (5) The bank's bonus structure for relationship managers is based 100 percent on revenue with no risk adjustment, (6) The bank's Board has not reviewed any AML metrics for the private banking division in 2 years.}
		
		\vspace{0.02cm}
		\textbf{What does this indicate about the bank's culture of compliance?}
		\vspace{0.02cm}
		
		\begin{itemize}
			\item A) Strong culture; compliance recommendations are considered
			\item B) Weak culture of compliance; revenue overrides AML; senior management overrides compliance; no Board oversight; compensation structure encourages risk-taking
			\item C) Only the bonus structure is problematic
			\item D) Only the lack of Board review is problematic
		\end{itemize}
		
		\vspace{0.02cm}
		\trap{Senior management can override compliance for commercial reasons.}
		
		\vspace{0.02cm}
		\correct{Answer: B - This is a textbook failure of compliance culture. Senior management overriding compliance for revenue is a critical failure. Board not reviewing AML metrics demonstrates lack of tone from the top. Compensation without risk adjustment encourages misconduct. Regulator would take enforcement action.}
		
		\vspace{0.02cm}
		\redflag{Senior management overriding compliance = culture of compliance failure.}
	\end{frame}
	
	% ============================================================
	% QUESTION 28 - ESTATE PLANNING FOR HNW CLIENTS
	% ============================================================
	
	\begin{frame}{Private Banking: Estate Planning Risks}
		\tiny
		\scenario{A private banking client establishes a trust for estate planning purposes with the following structure: Settlor: Client (age 75). Trustees: Client and his son. Beneficiaries: Client (life interest), then son and grandchildren. The client funds the trust with \$40 million in assets. Three months later, the client resigns as trustee, leaving son as sole trustee. Six months later, the trust sells \$30 million in assets and distributes the proceeds to the son personally (not as trustee). The son uses the funds to purchase a \$25 million yacht and \$5 million in cryptocurrency. The son then moves to a country with no extradition treaty. The client (father) claims he had no knowledge of the distributions. The bank's KYC files show the son has no independent source of wealth.}
		
		\vspace{0.02cm}
		\textbf{What red flags indicate that this trust structure is being used to move assets improperly?}
		\vspace{0.02cm}
		
		\begin{itemize}
			\item A) No red flags; estate planning is legitimate
			\item B) Rapid resignation of settlor as trustee; distributions to son personally (not as trustee); son has no source of wealth; son's relocation to non-extradition country; father's claim of no knowledge
			\item C) Only the yacht purchase is a red flag
			\item D) Only the cryptocurrency purchase is a red flag
		\end{itemize}
		
		\vspace{0.02cm}
		\trap{Trust distributions to beneficiaries are normal.}
		
		\vspace{0.02cm}
		\correct{Answer: B - Red flags: settlor rapid resignation, distributions to son in personal capacity (not as trustee), son's lack of independent wealth, relocation to non-extradition country. This suggests potential asset dissipation or fraud. Bank should file SAR.}
		
		\vspace{0.02cm}
		\redflag{Trust distributions to individual not in trustee capacity = potential fraud or money movement.}
	\end{frame}
	
	% ============================================================
	% QUESTION 29 - LEGACY AND SUCCESSION PLANNING
	% ============================================================
	
	\begin{frame}{Private Banking: Legacy and Succession Planning}
		\tiny
		\scenario{A private banking client is age 80 with a net worth of \$250 million. The client has three adult children. The client creates the following succession plan: (1) A trust for Child A with \$80 million, (2) A trust for Child B with \$80 million, (3) A trust for Child C with \$80 million, (4) A foundation with \$10 million for charitable purposes. All trusts are established in the Cook Islands. The trustees for all three trusts are the same law firm in the Cook Islands. The client gives each child a "letter of wishes" but refuses to provide copies to the bank. Six months after the client's death, the bank notices: (1) All three trusts have transferred their assets to new trusts in the Marshall Islands, (2) The new trusts have different trustees (company owned by Child A), (3) Child B has moved to a country with no financial disclosure laws, (4) Child C has purchased \$30 million in cryptocurrency through unregulated VASPs, (5) The foundation's \$10 million was wired to an account in the Seychelles with no charitable activity.}
		
		\vspace{0.02cm}
		\textbf{What succession planning red flags indicate potential money laundering or asset dissipation?}
		\vspace{0.02cm}
		
		\begin{itemize}
			\item A) No red flags; HNW clients can structure succession as they wish
			\item B) Migration to secrecy jurisdictions (Cook Islands to Marshall Islands); self-dealing (Child A as trustee); relocation to non-disclosure country; crypto purchases through unregulated VASPs; foundation funds diverted from charitable purpose
			\item C) Only the foundation diversion is a red flag
			\item D) Only the cryptocurrency purchases are a red flag
		\end{itemize}
		
		\vspace{0.02cm}
		\trap{Succession planning is private; no bank oversight.}
		
		\vspace{0.02cm}
		\correct{Answer: B - Multiple red flags: movement to increasingly secretive jurisdictions, self-dealing (beneficiary becomes trustee), relocation to non-disclosure jurisdiction, crypto through unregulated VASPs, charitable funds diverted to secrecy jurisdiction. This suggests possible asset hiding or money laundering.}
		
		\vspace{0.02cm}
		\redflag{Foundation funds diverted to secrecy jurisdiction = loss of charitable purpose + potential ML.}
	\end{frame}
	
	% ============================================================
	% QUESTION 30 - OFFSHORE TRUST NUANCES
	% ============================================================
	
	\begin{frame}{Private Banking: Offshore Trust Nuances}
		\tiny
		\scenario{A private bank is reviewing an offshore trust established in the Bahamas. The trust has the following features: (1) The settlor is a resident of Country X (no AML controls), (2) The trustee is a Bahamas trust company that is an affiliate of the bank, (3) The protector is a Swiss lawyer, (4) The beneficiaries are a class of "family members" with no names provided, (5) The trust has a "flee clause" allowing assets to be moved to another jurisdiction if the trust is threatened, (6) The trust has a "letter of wishes" that the settlor updates annually, (7) The trust's assets include \$50 million in cash, \$20 million in art stored in a Geneva freeport, and \$30 million in crypto held by an unregulated VASP. The trust refuses to provide the letter of wishes or identify any beneficiaries. The bank's relationship manager notes that the settlor has created and dissolved 6 similar trusts in the past 10 years, each lasting 18-24 months.}
		
		\vspace{0.02cm}
		\textbf{What features of this offshore trust present the highest money laundering risk?}
		\vspace{0.02cm}
		
		\begin{itemize}
			\item A) No features present risk; offshore trusts are legitimate
			\item B) Anonymous beneficiaries; flee clause (asset mobility); letter of wishes (secret instructions); rapid creation/dissolution of trusts; art in freeport (opaque asset); crypto with unregulated VASP; settlor from no-AML jurisdiction
			\item C) Only the crypto assets are a risk
			\item D) Only the freeport art is a risk
		\end{itemize}
		
		\vspace{0.02cm}
		\trap{Offshore trusts are standard for HNW international families.}
		
		\vspace{0.02cm}
		\correct{Answer: B - Highest risk features: (1) Anonymous beneficiaries (no CDD possible), (2) Flee clause allows rapid asset movement to avoid freezing, (3) Letter of wishes provides secret instructions not subject to review, (4) Pattern of trust creation/dissolution (each lasting 18-24 months suggests moving assets to avoid detection), (5) Freeport art (opaque valuation and ownership), (6) Crypto through unregulated VASP. This is a high-risk structure requiring EDD and likely SAR.}
		
		\vspace{0.02cm}
		\redflag{Flee clauses + anonymous beneficiaries + pattern of rapid trust dissolution = high risk.}
	\end{frame}
	
	% ============================================================
	% ANSWER KEY SUMMARY
	% ============================================================
	
	\begin{frame}{Answer Key Summary}
		\begin{multicols}{2}
			\tiny
			\begin{itemize}
				\item Q1: \correct{B}
				\item Q2: \correct{B}
				\item Q3: \correct{B}
				\item Q4: \correct{B}
				\item Q5: \correct{B}
				\item Q6: \correct{B}
				\item Q7: \correct{B}
				\item Q8: \correct{B}
				\item Q9: \correct{B}
				\item Q10: \correct{B}
				\item Q11: \correct{B}
				\item Q12: \correct{B}
				\item Q13: \correct{B}
				\item Q14: \correct{B}
				\item Q15: \correct{B}
				\item Q16: \correct{B}
				\item Q17: \correct{B}
				\item Q18: \correct{B}
				\item Q19: \correct{B}
				\item Q20: \correct{B}
				\item Q21: \correct{B}
				\item Q22: \correct{B}
				\item Q23: \correct{B}
				\item Q24: \correct{B}
				\item Q25: \correct{B}
				\item Q26: \correct{B}
				\item Q27: \correct{B}
				\item Q28: \correct{B}
				\item Q29: \correct{B}
				\item Q30: \correct{B}
			\end{itemize}
		\end{multicols}
		\vspace{0.05cm}
		\centering
		\greennote{30 Hard Questions - High Net Worth Individuals - Private Banking - Trusts - PEPs}
	\end{frame}
	
	% ============================================================
	% KEY CONCEPTS RECAP
	% ============================================================
	
	\begin{frame}{Key Concepts Recap - HNW Private Banking}
		\tiny
		\begin{itemize}
			\item \keyterm{Conflict of Interest:} RMs compensated on AUM may overlook red flags
			\item \keyterm{Source of Wealth vs Source of Funds:} Both require independent verification
			\item \keyterm{Trusts:} Identify settlor, all trustees, protector, ALL beneficiaries
			\item \keyterm{PEP Definition:} Includes foreign, domestic, international org PEPs, plus immediate family and close associates
			\item \keyterm{Foreign PEPs:} Always require EDD regardless of time since leaving office
			\item \keyterm{Domestic PEPs:} Risk-based assessment; red flags require EDD
			\item \keyterm{Offshore Centers:} Red flags: complex structures, round tripping, rapid transfers, shell companies
			\item \keyterm{High Value Assets:} Rapid resale at loss = money laundering (loss is cost)
			\item \keyterm{Insurance ML:} Policy loan layering; early surrender for clean funds
			\item \keyterm{Secured Loans:} Illicit collateral + loan proceeds = clean funds
			\item \keyterm{Correspondent Banking:} Must identify underlying private banking clients; unauthorized nesting = terminate
			\item \keyterm{Culture of Compliance:} Senior management cannot override compliance for revenue
			\item \keyterm{Trust Red Flags:} Flee clauses, anonymous beneficiaries, rapid creation/dissolution, self-dealing
			\item \keyterm{Freeports:} Art storage with no physical possession = opacity risk
			\item \keyterm{Family Offices:} Unidentified third parties, multi-jurisdiction structures = EDD required
		\end{itemize}
		\vspace{0.05cm}
		\centering
		\greennote{Good luck on your CAMS exam!}
	\end{frame}
	
\end{document}